We now take a further step in weakening the conditions on solutions of the Dirichlet problem: we state the problem in variational form. This is the same process which was used in discussing weak solutions of conservation laws in Chapter~\ref{chap:conservation-laws-shocks}. The first step is to create a bilinear form from the differential operator $L$ using integration by parts. Let $\phi\in\mathcal{D}(\Omega)$ and $u\in H^{2k}(\Omega)$, then
\begin{equation}\label{eq:dirichlet-variational-integration-by-parts}
\begin{aligned}
\int_\Omega\phi Lu\,\mathrm{d}\bm{x}
&=\sum_{0\leq|\sigma|,|\gamma|\leq k}(-1)^{|\sigma|}\int_\Omega\phi D^\sigma(a_{\sigma\gamma}(\bm{x})D^\gamma u)\,\mathrm{d}\bm{x}\\
&=\sum_{0\leq|\sigma|,|\gamma|\leq k}\int_\Omega a_{\sigma\gamma}(\bm{x})D^\gamma uD^\sigma\phi\,\mathrm{d}\bm{x}.
\end{aligned}
\end{equation}


With this in mind we define
\begin{equation}\label{eq:elliptic-dirichlet-bilinear-form-definition}
B[v,u]:=\sum_{0\leq|\sigma|,|\gamma|\leq k}\int_\Omega a_{\sigma\gamma}(\bm{x})D^\gamma uD^\sigma v\,\mathrm{d}\bm{x}
\end{equation}
to be the \textbf{bilinear form} associated with the elliptic partial differential operator $L$. Note that $B[v,u]$ is well defined for $u$ and $v$ that are merely in $H^k(\Omega)$. With this in mind, we give the following definition of yet another type of solution of the Dirichlet problem.

\begin{Definition}
\label{def:weak-dirichlet-solution}
Let $\Omega\subset\mathbb{R}^n$ be a bounded domain and suppose $f\in H^{-k}(\Omega)$ is given. A function\[u\in H_0^k(\Omega)\]is a \textbf{weak solution} of the Dirichlet problem if
\begin{equation}\label{eq:weak-dirichlet-variational-equation}
B[v,u]=f(v)
\end{equation}
for every $v\in H_0^k(\Omega)$.
\end{Definition}

\begin{Remark}
\label{rem:complex-dirichlet-sesquilinear-extension}
We can extend the bilinear form $B$ on the real Hilbert space $H_0^k(\Omega)$ to be a \emph{sesquilinear form} on the complex Hilbert space (also denoted $H_0^k(\Omega)$) by letting
\begin{equation}\label{eq:complex-elliptic-dirichlet-sesquilinear-form}
B[v,u]:=\sum_{0\leq|\sigma|,|\gamma|\leq k}\int_\Omega a_{\sigma\gamma}(\bm{x})D^\gamma u\overline{D^\sigma v}\,\mathrm{d}\bm{x}.
\end{equation}
We will be rather sloppy about the distinction between complex-valued and real-valued functions, using the same notation for the spaces and the bilinear forms. Since we are mainly interested in discussing real solutions to partial differential equations, we will only go to complex spaces when forced to, such as when using Fourier transforms or discussing spectral theory.
\end{Remark}

Note that by using the calculations of \eqref{eq:dirichlet-variational-integration-by-parts} (though this time with a function $v\in H_0^k(\Omega)$ in place of $\phi\in\mathcal{D}(\Omega)$) we can see that any strong solution of the Dirichlet problem is automatically a weak solution. However, since we require so much less smoothness of weak solutions than strong ones, it will be far easier to show that weak solutions exist. Once we have done this, we will be able to show that if $\Omega$, $f$ and the coefficients $a_{\sigma\gamma}$ are sufficiently ``nice,\textquotesingle\textquotesingle\ the weak solution is, in fact, a strong solution or a classical solution.
